\chapter{Introducere}
\label{cap:cap1}

Introducerea va avea aproximativ 2 (două) pagini care vor conține motivația alegerii temei, relevanța și contextul temei alese, obiectivele generale ale lucrării, metodologia și instrumentele utilizate și o scurtă descriere a structurii lucrării (titlul capitolelor, o scurtă descriere și legătura dintre acestea).

În acest capitol nu se introduc figuri, table sau listing-uri de cod. Pot fi în schimb referite!

%\sout{text de șters}

\section{Recomandări generale pentru redactare}

\begin{description}
    \item[Font:]  Times New Roman, mărimea 12
    \item[Spațiere:] 1 între rânduri
    \item[Margini:] 2 cm (sus, jos, stânga, dreapta)
    \item[Număr de pagini:] Maxim 30 pagini fără rezumat, bibliografie și anexe
    \item[Stil bibliografic:] IEEE
    \item[Format:] Word sau LaTeX (în funcție de preferințe)
\end{description}

\section{Structura lucrării scrise}

\begin{itemize}
    \item Rezumatul lucrării (max. 1 pg.)
    \item 1. Introducere (aprox. 2 pg.)
    \item 2. Fundamentarea teoretică și soluții similare (aprox. 5 pg.)
    \item 3. Soluția propusă (aprox 14 pg.)
    \item 4. Testarea soluției și rezultate experimentale (aprox. 7 pg.)
    \item 5. Concluzii (max. 2 pg.)
    \item Bibliografie
    \item Anexe
\end{itemize}

\section{Recomandări conținut introducere}

\begin{itemize}
    \item Contextul și importanța temei alese.
    \item Obiectivele generale și specifice ale lucrării.
    \item Metodologia de cercetare utilizată.
    \item Structura generală a lucrării.
\end{itemize}